%!TEX program = xelatex
\documentclass[12pt, b5paper]{article}
\usepackage{ctex}

\usepackage{geometry}
\geometry{b5paper,left=2cm,right=1cm,top=2.5cm,bottom=2.5cm}
\usepackage{chemformula} %化学公式宏包
\usepackage[version=4]{mhchem} %化学公式宏包
\usepackage{siunitx} %数字，单位宏包
\sisetup{
	separate-uncertainty = true,
	inter-unit-product = \ensuremath{{}\cdot{}}
} %单位间加小圆点
\usepackage{tikz} %Tikz绘图包
\usetikzlibrary{cd}
\usetikzlibrary{chains}
\usetikzlibrary{arrows}
\usetikzlibrary{arrows.meta} %画箭头
\usetikzlibrary{commutative-diagrams} %交换图
\usetikzlibrary{positioning,automata}
\usepackage[all]{xy}
\usepackage{booktabs} %三线表
\usepackage{multirow} %合并单元格
\usepackage{makecell} %表格内强制换行
\usepackage{array}
\usepackage{booktabs} %控制表格行高
\usepackage{colortbl}  %彩色表格需要加载的宏包
\usepackage{amsmath}
\usepackage{unicode-math}
\usepackage{fontspec} %字体
\newcommand*{\cred}{\textcolor{red}}

\setmainfont{Times New Roman}
\setmathfont{XITS Math}

\setlength{\parindent}{0pt} %无缩进
\usepackage{setspace}
\usepackage{fancyhdr} %设置页码
\usepackage{lastpage} %获取总页数
\pagestyle{fancy} %fancyhdr宏包新增的页面风格
\renewcommand\headrulewidth{0pt} %取消页眉中的装饰分割线
\fancyhf{}
\cfoot{\footnotesize 第 \thepage\ 页(共 \pageref{LastPage}页)}%当前页 of 总页数

\begin{document}
\begin{spacing}{1.625}

\centerline{{\huge 河北农业大学课程考试试卷}}

\centerline{\large 2022-2023学年第\underline{\ 1\ }学期}

考试科目：\underline{物理化学与胶体化学} \quad 姓名：\underline{ \qquad \qquad} \quad 学号：\underline{ \qquad \qquad \qquad}

学院：\underline{理工学院} \quad 专业班级：\underline{\qquad \qquad \qquad \qquad} \quad 卷别：\cred{A} 考核方式：\cred{开卷}

\bigskip


1. 简答题(10分)

有人认为“空调可以从低温热源吸热放给高温热源，这与克劳修斯表述不符”。请指出错误所在。

\bigskip

\bigskip

2. 简答题(10分)

电解质溶液的电导测定有何实际意义？

\bigskip

\bigskip

3. 简答题(10分)

在常温下，对于NaCl水溶液，考虑电离和不考虑电离时，体系的独立组分数和物种数的区别。

\bigskip

\bigskip

4. 简答题(10分)

合成氨的反应是一个放热反应，降低反应温度有利于提高转化率。但实际生产中这一反应都是在高温高压和有催化剂存在的条件下的进行的，为什么？

\bigskip

\bigskip

5. 简答题(10分)

已知水的三相点为273.16K，611Pa。水的升华热$\Delta_\mathrm{sub}H_\mathrm{m}=\SI{51.06}{kJ.mol^{-1}}$。某日早晨气温-10℃，当大气中的水蒸气分压为300Pa时，霜能否升华变为水气？

\newpage

6. 主观题(20分)

乙烷脱氢生产乙烯 \ch{C2H6 (g) -> C2H4 (g) + H2 (g)}

(1) 判断该反应在常温常压下能否自发进行。

(2) 判断该反应在500K下能否自发进行。相关热力学数据查课本附录。

\bigskip

\bigskip

7. 主观题(15分)

银器长期放置，光泽暗淡，表面会发黄、变黑，是由于银长期露置空气中，被氧化所致。此类情况在各种博物馆展出的古银器上尤为常见。

\centerline{\ch{4 Ag + O2 -> 2 Ag2O}}

将银的氧化反应设计成电池，(1) 正负极应分别选用什么电极？(2) 写出电池的表示式和电极反应。(3) 已知$\varphi_\mathrm{AgO/Ag}^\ominus = \SI{0.34}{V}$，计算该电池的标准电动势。

\bigskip

\bigskip

8. 主观题(15分)

试计算在珠穆朗玛峰顶上的沸水中煮熟一个鸡蛋需要多长时间。

己知组成蛋的蛋白质的热变作用为一级反应，其活化能约为\SI{85}{kJ.mol^{-1}}，在$p^\ominus$下的沸水中煮熟一个鸡蛋需要\SI{10}{min}。水的摩尔蒸发热$\Delta H_\mathrm{m} = \SI{41.05}{kJ.mol^{-1}}$。珠穆朗玛峰顶大气压力$p=\SI{31}{kPa}$。



\end{spacing}
\end{document}
